\documentclass[11pt]{letter}

\usepackage[T1]{fontenc}
\usepackage[margin=1in]{geometry}

\signature{Rylan Malarchick\\Department of Engineering Physics\\Embry-Riddle Aeronautical University\\Daytona Beach, FL 32114\\malarchr@my.erau.edu}

\begin{document}

\begin{letter}{Editors\\Journal of Computational Physics}

\opening{Dear Editors,}

Please consider the enclosed manuscript, ``Hardware-Aware Performance
Characterization of Small Dense Lindblad Propagation for Near-Term Quantum
Control,'' for publication in the \textit{Journal of Computational Physics}.

This paper is a computational physics comparison study of dense Lindblad
propagation in the small-matrix regime most common in near-term transmon
control. The main result is not a single hardware winner, but a measured map
of operating regimes across two x86 CPUs, two NVIDIA GPUs, and a Tang Nano
20K FPGA. The manuscript combines working-set analysis with empirical
Roofline context, separates host-visible and kernel-only GPU timing, and
includes an end-to-end GRAPE-style benchmark against released QuTiP 5.2.3.
That application benchmark matters because it shows where kernel speedups do
and do not survive into control workflows: the optimized C path wins at
$d=3$ and $d=9$, but loses at $d=27$ because propagator construction, not
propagation, dominates the landscape-point cost.

I believe the manuscript fits JCP because it is centered on computational
methods, performance characterization, and comparison against established
approaches. The paper is intended to be useful to authors of quantum-control
software who need to know which hardware helps in this specific computational
regime, which timing model is relevant, and where further algorithmic work is
still needed.

This manuscript is original, has not been published previously, and is not
under consideration for publication elsewhere. All results reported in the
manuscript are traceable to code, timing logs, and figure-generation scripts
available in the accompanying repository. A declaration of generative AI and
AI-assisted technologies is included in the manuscript in accordance with the
journal's policy. I confirm that I have no competing financial or
non-financial interests related to this submission.

Thank you for your consideration.

\closing{Sincerely,}

\end{letter}

\end{document}
